\documentclass[11pt]{article}
\usepackage[margin=1in]{geometry}
\usepackage{amsmath,amssymb}
\usepackage{booktabs}
\usepackage{enumitem}
\usepackage{hyperref}
\usepackage{xcolor}

\newcommand{\perc}{\mathrm{perc}}
\newcommand{\self}{\mathrm{self}}

\title{Cross-Network Effects Between Perceived Slices and Self-Reported Network\\[4pt]
\large Work Log --- April 13, 2026}
\author{Jinwoo Cho \and Nynke Niezink}
\date{}

\begin{document}
\maketitle

%% ==========================================================
\section{Overview}
%% ==========================================================

We implemented cross-network effects between perceived slices $Y^{(p)}$ and
the self-reported network $Y^{(\self)}$ by reusing RSiena's existing
\texttt{crprod} infrastructure (the \texttt{nonSymmetricNonSymmetricObjective}
effect group), avoiding any C++ changes.

\paragraph{Notation.}
The CSS tensor is $x_{pjk}$ where $p$ = perceiver, $j$ = sender, $k$ = receiver.
The perceived slice for perceiver $p$ is the $n \times n$ one-mode network
$Y^{(p)}_{jk} = x_{pjk}$, and the self-reported network is
$Y^{(\self)}_{jk} = x_{jjk}$.
In the standard RSiena two-network notation (manual Section~12.1.3),
the dependent network is $X = Y^{(p)}$ and the explanatory network is
$W = Y^{(\self)}$.
The focal actor making a ministep in $Y^{(p)}$ is sender $j$ ($j \neq p$,
since row $p$ is structurally fixed).

%% ==========================================================
\section{Dyadic Cross-Network Effects}
%% ==========================================================

\subsection{Entrainment / crprod (\texttt{crprod})}

\textbf{RSiena formula} (manual effect 1):
\[
  s_{j}^{\,\mathrm{net}}(x) \;=\; \sum_{k} x_{jk}\, w_{jk}.
\]
\textbf{Three-way form}:
\[
  s_{j}^{(p)}(x) \;=\; \sum_{k} x_{pjk} \;\cdot\; x_{jjk}.
\]
\textbf{Interpretation.}
Perceptual accuracy: how much perceiver~$p$'s view of $j$'s outgoing ties
matches $j$'s own self-reported ties.
A positive parameter means the perceived network tends to align with
the self-report.

% ----------------------------------------------------------
\subsection{Reciprocal entrainment (\texttt{crprodRecip})}

\textbf{RSiena formula} (manual effect 2):
\[
  s_{j}^{\,\mathrm{net}}(x) \;=\; \sum_{k} x_{jk}\, w_{kj}.
\]
\textbf{Three-way form}:
\[
  s_{j}^{(p)}(x) \;=\; \sum_{k} x_{pjk} \;\cdot\; x_{kkj}.
\]
\textbf{Interpretation.}
Generalised exchange:
the existence of $k$'s self-reported tie \emph{to}~$j$
(i.e.\ $k$ says ``I have a tie to $j$'')
promotes perceiver~$p$ perceiving that $j$ has a tie \emph{to}~$k$.
A positive parameter indicates that being named by $k$ makes it more
likely that $j$ is perceived as sending a tie to $k$.

% ----------------------------------------------------------
\subsection{Mutual entrainment (\texttt{crprodMutual})}

\textbf{RSiena formula} (manual effect 3):
\[
  s_{j}^{\,\mathrm{net}}(x) \;=\; \sum_{k} x_{jk}\, w_{jk}\, w_{kj}.
\]
\textbf{Three-way form}:
\[
  s_{j}^{(p)}(x) \;=\; \sum_{k} x_{pjk} \;\cdot\; x_{jjk} \;\cdot\; x_{kkj}.
\]
\textbf{Interpretation.}
Accuracy conditional on mutual self-report:
the perceived tie $j \to k$ aligns with reality only when both $j$ and $k$
self-report ties to each other.
A positive parameter suggests that mutual self-reported ties are
especially visible to perceivers.

% ----------------------------------------------------------
\subsection{Entrainment $\times$ in-degree activity (\texttt{crprodInActIntn})}

\textbf{RSiena formula} (manual effect 13):
\[
  s_{j}^{\,\mathrm{net}}(x) \;=\;
    \sum_{k} x_{jk}\, w_{jk}\,
    \bigl(\sqrt{w_{+j}} - \sqrt{\bar{w}}\bigr),
    \quad p = 2.
\]
\textbf{Three-way form}:
\[
  s_{j}^{(p)}(x) \;=\;
    \sum_{k} x_{pjk} \;\cdot\; x_{jjk} \;\cdot\;
    \Bigl(\sqrt{\textstyle\sum_{\ell} x_{\ell\ell j}}
          - \sqrt{\bar{w}}\Bigr),
\]
where $\sum_{\ell} x_{\ell\ell j}$ is the self-reported in-degree of $j$
(how many others report a tie to $j$), and $\bar{w}$ is the mean degree
of $Y^{(\self)}$.

\textbf{Interpretation.}
Moderation of perceptual accuracy by self-reported popularity:
the alignment between perceived and self-reported ties is
stronger for actors who are popular (high in-degree) in the
self-reported network.

%% ==========================================================
\section{Degree-Related Cross-Network Effects}
%% ==========================================================

The following effects involve degree statistics of $W = Y^{(\self)}$.
In all cases, $W$-degrees are centered by the average degree $\bar{w}$ of
$Y^{(\self)}$.

\subsection{$W$ in-degree on $X$-popularity (\texttt{inPopIntn})}

\[
  s_{j}^{(p)}(x) = \sum_{k} x_{pjk}
    \bigl(\sqrt{w_{+k}} - \sqrt{\bar{w}}\bigr),
  \quad
  w_{+k} = \sum_{\ell} x_{\ell\ell k}.
\]
Actors who are popular in the self-reported network attract more perceived
incoming ties.

\subsection{$W$ in-degree on $X$-activity (\texttt{inActIntn})}

\[
  s_{j}^{(p)}(x) = x_{pj+}
    \bigl(\sqrt{w_{+j}} - \sqrt{\bar{w}}\bigr),
  \quad
  w_{+j} = \sum_{\ell} x_{\ell\ell j}.
\]
Actors who are popular in the self-reported network are perceived as more
active (higher perceived out-degree).

\subsection{$W$ out-degree on $X$-popularity (\texttt{outPopIntn})}

\[
  s_{j}^{(p)}(x) = \sum_{k} x_{pjk}
    \bigl(\sqrt{w_{k+}} - \sqrt{\bar{w}}\bigr),
  \quad
  w_{k+} = \sum_{\ell} x_{kk\ell}.
\]
Perceived ties tend toward actors who are active (high out-degree)
in the self-reported network.

\subsection{$W$ out-degree on $X$-activity (\texttt{outActIntn})}

\[
  s_{j}^{(p)}(x) = x_{pj+}
    \bigl(\sqrt{w_{j+}} - \sqrt{\bar{w}}\bigr),
  \quad
  w_{j+} = \sum_{\ell} x_{jj\ell}.
\]
Actors who are active in the self-reported network are perceived as
more active.

%% ==========================================================
\section{Mixed Triadic Effects}
%% ==========================================================

These effects capture transitive-like configurations across the
perceived and self-reported networks.
Below, $X = Y^{(p)}$ and $W = Y^{(\self)}$.

\subsection{From $W$ agreement (\texttt{from})}

\[
  s_{j}^{(p)}(x) = \sum_{h \neq k} x_{pjk}\, w_{jh}\, w_{kh}
  = \sum_{h \neq k} x_{pjk}\, x_{jjh}\, x_{kkh}.
\]
Structural equivalence in outgoing self-reports:
perceived tie $j \to k$ is promoted when $j$ and $k$ self-report ties to
the same third actors $h$ (i.e.\ $j$ and $k$ agree on who to nominate).

\subsection{$W$ to agreement (\texttt{to})}

\[
  s_{j}^{(p)}(x) = \sum_{h \neq k} x_{pjk}\, w_{hj}\, x_{phk}
  = \sum_{h \neq k} x_{pjk}\, x_{hhj}\, x_{phk}.
\]
Mixed $W$--$X$ two-path closure:
if there is a self-reported tie $h \to j$ and a perceived tie $h \to k$
by perceiver $p$, this promotes the perceived tie $j \to k$.

\subsection{Closure of $W$ (\texttt{closure})}

\[
  s_{j}^{(p)}(x) = \sum_{h \neq k} x_{pjk}\, w_{jh}\, w_{hk}
  = \sum_{h \neq k} x_{pjk}\, x_{jjh}\, x_{hhk}.
\]
Closure of $W$--$W$ two-paths:
if $j$ self-reports a tie to $h$ and $h$ self-reports a tie to $k$,
this promotes perceiver~$p$ perceiving a tie $j \to k$.
``Self-reported two-paths tend to become perceived ties.''

\subsection{XWX closure (\texttt{cl.XWX})}

\[
  s_{j}^{(p)}(x) = \sum_{h \neq k} x_{pjk}\, x_{pjh}\, w_{hk}
  = \sum_{h \neq k} x_{pjk}\, x_{pjh}\, x_{hhk}.
\]
If perceiver~$p$ sees $j \to h$ and $h$ self-reports a tie to $k$,
this promotes perceiver~$p$ also perceiving $j \to k$.

%% ==========================================================
\section{Implementation Notes}
%% ==========================================================

\begin{enumerate}[leftmargin=*]
\item
  All effects are generated automatically by calling
  \texttt{createEffects("nonSymmetricNonSymmetricObjective",
  selfName, name = sliceName)} inside \texttt{threeWayNet()} in
  \texttt{effects.r}.

\item
  The shared parameter framework applies:
  $Y[\text{shared}]$ is the canonical parameter holder
  (\texttt{sharedDup = FALSE}), and
  $Y[2], \dots, Y[K]$ are hidden duplicates
  (\texttt{sharedDup = TRUE}).

\item
  Usage example:
  \begin{verbatim}
  eff <- includeEffects(eff, crprod,
                        interaction1 = "Y[self]",
                        name = "Y[shared]")
  \end{verbatim}

\item
  No C++ modifications were needed.
  The C++ side resolves
  \texttt{pState->pNetwork("Y[self]")} by name string, and
  \texttt{initializeFRAN.r} already registers $Y^{(\self)}$ as a
  separate one-mode network.
\end{enumerate}

%% ==========================================================
\section{Bug Fix: Structural Zeros in Perceived Slices}
%% ==========================================================

We identified a bias in the parameter recovery simulation
(true $\beta = 0.6$, estimated $\hat\beta \approx 0.4$).
The cause was that the structural-10 marking of perceiver~$p$'s own row
in slice $Y^{(p)}$ was applied only in \texttt{effects.r}
(for effect-table generation) but \emph{not} in \texttt{initializeFRAN.r}
(for the actual C++ data).

Without structural zeros on row $p$ of $Y^{(p)}$:

\begin{itemize}[leftmargin=*]
\item Actor $p$ could make ministeps in $Y^{(p)}$, modifying $x_{ppk}$.
\item The same cells $x_{ppk}$ also live in $Y^{(\self)}$,
      causing double-counting.
\item The crprod statistic was inflated by the $j = p$ self-correlation
      term $\sum_k (x_{ppk})^2$, leading to downward bias in
      $\hat\beta_{\texttt{crprod}}$.
\end{itemize}

\noindent\textbf{Fix:} In \texttt{initializeFRAN.r},
\texttt{.expandThreewayToOneModes()}, after extracting
\texttt{arr3 <- depvar[k,{},{}]}, we now set
\texttt{arr3[k, , ] <- 10} (structural zero) for all time points.
The distance computation was also updated to exclude row $k$.

%% ==========================================================
\section{Bug Fix: DGP Reciprocity at the Structural-Row Boundary}
%% ==========================================================

After fixing the structural-zero issue above, a second bias appeared
in the $Y^{(\self)}$ parameters:
self out-degree bias $\approx -0.58$,
self reciprocity bias $\approx +0.35$
(across 10 replications).

\paragraph{Root cause.}
In the data-generating process (DGP), when perceiver~$p$ evaluates
toggling the perceived tie $x_{pjk}$ with $k = p$, the reciprocity
contribution is
\[
  \beta_{\mathrm{recip}} \cdot x_{p,\,k,\,j}
  \;=\; \beta_{\mathrm{recip}} \cdot x_{ppj},
\]
which is the \emph{actual} self-reported value.
However, in RSiena's C++ engine, the cell $Y^{(p)}_{pj}$ is
structural zero (row~$p$ of $Y^{(p)}$ is entirely structural).
RSiena therefore always evaluates the reciprocity contribution as~0
for $k = p$.

This mismatch means the DGP generates data under a model where
perceived ties $x_{pjp}$ are attracted toward self-reported ties
$x_{ppj}$ via reciprocity, but RSiena cannot capture this mechanism.
The spurious signal propagates through joint estimation (via crprod)
and biases the $Y^{(\self)}$ parameters.

\paragraph{Fix.}
In the DGP's \texttt{delta\_f\_perceived\_ijk()}, set the reciprocity
term to~0 when $k = p$:
\begin{verbatim}
recip_val <- if (k == i) 0 else x[i, k, j]
\end{verbatim}
This aligns the DGP with RSiena's structural-zero convention: the
perceived reciprocity statistic is blind to row~$p$ of $Y^{(p)}$.

%% ==========================================================
\section{Bug Fix: Name Map Misalignment in \texttt{initializeFRAN}}
%% ==========================================================

The $Y^{(\self)}$ parameter estimates appeared identical to perceived
parameter estimates (e.g., self outdegree~$\approx -1.63$ instead of
true~$-0.98$, matching perceived~$-1.68$).
The cause was a row-ordering bug in \texttt{initializeFRAN.r}.

\paragraph{Root cause.}
The \texttt{Y[shared]} $\to$ \texttt{Y[1]} name remap saved a name map
(\texttt{z\$\,.shared\_name\_map}) \emph{before} the
\texttt{split + rbind} reordering of \texttt{requestedEffects}.
The \texttt{split + rbind} groups effects by \texttt{name} and
reorders them by \texttt{depvarnames}, which moves cross-network
effects (crprod) from the end of the table into each slice's group.
This changes the row order, but the saved map still reflected the
\emph{old} order.

When \texttt{terminateFRAN} restored the names from the saved map,
the misaligned mapping caused $Y^{(\self)}$'s effects to receive
names from perceived slice~16 (or whichever slice happened to occupy
the corresponding position). The theta values were therefore correct
for the position, but the labels were wrong---so
$Y^{(\self)}$'s true theta was plotted under a perceived name,
and a perceived theta was plotted under the $Y^{(\self)}$ label.

\paragraph{Fix.}
Save the name map \emph{after} the \texttt{split + rbind} by applying
the reverse remap (\texttt{Y[1]} $\to$ \texttt{Y[shared]}) to the
already-reordered \texttt{requestedEffects\$name}.

\end{document}
